

\section{MMU-based Page Access Obliviousness}\label{sec:design:mmu}

\key{
	As the foremost layer of defense, \thename\ thwarts controlled-channel attacks by providing \gls{pao}~\cite{pao}, ensuring that page-level memory access traces generated obtained through controlled-channel attacks are indistinguishable.
}
This are achieved by enforcing identical sequences of memory traces or obfuscating memory traces based on random page accesses.

\key{
	\thename\ efficiently combines both tasks by leveraging the huge pages (\textbf{O2}) and direct access to the \gls{mmu} (\textbf{O1}).
}
With \gls{pao} in place, \thename\ not only renders the profiling phase of attacks on \gls{cvm} (\cref{sec:overview:threat}) infeasible but also stops in-track a plethora of attacks~\cite{} that depends on the controlled channel as core attack component.


% Mreover, unlike compiler-based instrumentation, our Unikernel design fully leverage textbf{O1} to make it efficient.
% \thename{} leverage the huge mappings (\textbf{O2}) to render the paging scheme very efficient.



% Implementing a paging-based \gls{pao} scheme requires \thename\ to tackle three challenges.
% The first paging scheme must include code and data pages and the OS page table.
% Second, it must secure the location of trap and interrupt handlers (i.e., the \gls{isr}) from fingerprinting attacks, given that their mappings must always be present.
% Finally, \thename{} must also retrofit the integrity of huge mappings in the existing \gls{cvm} architectures that allow the hypervisor to change memory mapping sizes arbitrarily.
% In the next section, we will go into details \thename's solution to these problems.

\begin{figure}[t]
	\centering
	\includegraphics[width=.95\columnwidth]{figures/incongitos/concept.pdf}
	\caption{Objective of \thename's paging scheme}
	\label{fig:concept}
\end{figure}

% \key{
% 	Since hardware interrupts are virtualized by the hypervisor, \thename must ensure that the hypervisor sends interrupts promptly.
% }
% \thename uses the secure timer (Secure \gls{tsc}) feature provided by \gls{sevsnp} to reliably obtain how much time has passed since the last validation.
% % Fortunately, SEV-SNP provides a secure timer (Secure TSC) mechanism that the hypervisor cannot manipulate. 
% Once Secure TSC is enabled, \texttt{rdtsc} computes the TSC with values outside the hypervisor's control~\cite{amd64}.
% \thename then employs a threshold $\Delta$ as the maximum time between each validation.
% When the time measured since the last validation exceeds $\Delta$, \thename also terminates the execution.


\subsection{Efficient ORAM-backed paging}
% \hdr{Whole-system ORAM-backed paging}

% \key{
% 	\thename takes full advantage of the gracious 2MB physical space provided by \gls{apage} to limit the amount of costly \gls{oram}-backed page swaps during runtime.
% }

% By ensuring that all in-\gls{cvm} access via \glspl{gva} is confined to a corresponding \gls{gpa} within \gls{apage}, \thename achieves full-system \gls{pao}.
% In other words, all in-\gls{cvm} memory accesses must land on one of the 512 active subpages that are mapped into \gls{apage}'s \gls{gpa}.
% An exception to this rule is the page swapping operations; with the now 2MB-limited currently active subpages, a scheme for paging in/out of the subpages is indispensable for practicality.

\key{By exploiting direct MMU access, \thename decouples the paging granularity between the host and guest systems.
	From the untrusted hypervisor's perspective, \gls{apage} is backed by a single huge page, but from the guest's viewpoint, each guest page is managed at a 4KB granularity.}
This empowers the guest kernel to self-handle page faults at the 4KB level without detection by the host.
Physical addresses outside of \gls{apage} serve as \gls{oram}-backed memory, regardless of whether the pages are backed by 4KB or 2MB.

\thename then set up the OS page table, such that all \glspl{gva} that are not mapped to \gls{apage} would incur a \emph{page fault}.
Upon accessing the faulting \gls{gva}, \thename's page fault handler covertly transfers the page's contents from \gls{oram}-backed storage to a physical page within \gls{apage} using Path ORAM algorithm.
It then updates the page table to map the \gls{gva} to the corresponding \gls{gpa}.
\key{Once the \gls{gva} is mapped, the hardware \gls{mmu} automatically handles address translation, resulting in zero overhead compared to standard \gls{cvm} execution.}



\hdr{ORAM-backed page management}
\todo{Some background on ORAM}


\key{\thename's page table scheme retrofits the unused bits in its \gls{pte} to encode information for \thename ORAM-backed paging.}
\autoref{fig:oram-pte} shows \thename \gls{pte} encoding inside the page table.
A \gls{pte} that maps an ORAM-backed \gls{gva} may represent three different states.
When the present bit is set, then the mapping is active, i.e., the \gls{pte} points to a \gls{pfn} inside \gls{asa}.
A non-present PTE may be \emph{unallocated} or \emph{ORAM-backed}.
\thename leverages an unused bit in the \gls{pte} to indicate that the backing \gls{pfn} is residing inside the ORAM tree (i.e., ORAM-backed bit), and an ORAM access must be performed to fetch the page.
If the ORAM-backed bit is unset, \thename treats the \gls{pte} as unallocated and performs subpage allocation.
Subpage allocation performs subpage replacement to obtain an empty subpage within \gls{asa} without fetching a new subpage from the stash.
% \thename's paging scheme uses this bit to decide whether to allocate a new block or fetch a block from the ORAM tree with the Path ORAM algorithm.
% and an ORAM access must be performed to fetch the block. 

The Path ORAM algorithm maintains a PosMap structure that maps ORAM block addresses to their perspective leaf index in the ORAM tree (\cref{sec:bg:oram}).
Similarly, the ORAM scheme of \thename also uses a PosMap but minimizes the memory and performance overheads of PosMap management by integrating them into the inherent overheads of paging.
\key{
	Specifically, once a subpage is evicted from \gls{asa} (\CircledN{3} in \autoref{fig:ovma-fault}), its \gls{pte} is encoded as ORAM-backed, and its bits[51-12] encode the ORAM leaf index such that \thename can quickly retrieve the leaf index when \gls{pager} is invoked during the software page table walk without additional lookups.
}
In other words, the PosMap of ORAM is now replaced with the page table in \thename's ORAM scheme.


\hdr{Recursive-ORAM-backed software page walk}
Interestingly, the subpage access pattern of a software page table walk for handling a faulting ORAM-backed \gls{gva} resembles that of a recursive ORAM scheme.
Recursive ORAM algorithms break down the PosMap data structure into smaller blocks by a factor of $X$ and store them into another ORAM that uses a smaller PosMap.
This smaller PosMap is then, once again, broken down and stored into another ORAM.
The process is repeated recursively for $r$ times until the smallest PosMap is achieved, which is then stored inside secure memory.
% and store them into ORAMs, leading to progressively smaller PosMaps. 
To perform recursive ORAM access, PosMaps are accessed, starting from the smallest PosMap.
For PosMaps backed by ORAM, ORAM accesses are performed to retrieve the PosMap ORAM block that stores the leaf ID of the next ORAM.
A variation of this approach stores PosMap blocks alongside ORAM blocks within a \emph{unified} ORAM tree structure~\cite{freecursiveoram}, similar to \thename's single \gls{pagepool} both page tables and data pages.

Similarly to a Recursive ORAM access, \thename leverages ORAM accesses to obliviously retrieve the subpage containing the \gls{gpt} required for mapping the faulting \gls{gva}.
As we will discuss in \cref{sec:oram-encoding}, the \gls{gpt} now also functions as a PosMap with its retrofitted \gls{pte} scheme.
\key{
	{This means that \thename's software page table walk is in fact, a Recursive ORAM scheme with recursion level $r = 4$ levels of reduction factor $X = 512$ (the number of \gls{pte} per page table), where both data and PosMap ORAM blocks are 4KB and are stored in a unified ORAM tree.}
}





\subsection{Additional design considerations}
In this subsection, we discuss additional aspects of the design and implementation of \thename's MMU-assisted \gls{pao} scheme that were left out for a more concise presentation.

\hdr{Post-boot-time rerandomization}
After \thename's paging scheme is initialized, if the attacker knows of the initial sequence of page accesses, they may be able to deduce the initial physical memory layout.
This is because the page fault handler fetches the requested page into the code/data scratchpad at a specific offset, allowing the attacker to obtain the ciphertext update.
To address this issue, \thename\ rerandomize the scratchpads by randomizing the subpages and the respective \gls{gva}-to-\gls{gpa} mappings.
Combined with adaptive rerandomization, this ensures that it is extremely challenging for an attacker to infer the memory layout.



\hdr{Handling interrupt service routines}
Interrupt and trap handlers require their associated code and data pages to be mapped throughout kernel execution to handle hardware events.
As a result, these pages must be excluded from \thename's paging scheme to ensure correct kernel execution.
\thename\ is designed to allow certain kernel address regions to be excluded from its paging scheme.

It is important to note that since the hypervisor can arbitrarily inject virtual interrupts, the profiling of \glspl{isr} cannot be prevented through \gls{pao} alone.
However, as the remaining code and data are protected by \gls{pao}, the access pattern of the \glspl{isr} does not provide useful information to the attacker.

\hdr{Page replacement algorithm}
\key{
	\thename's subpage replacement scheme is based on CLOCK~\cite{} page replacement algorithm that approximates LRU with low costs to cope with frequent page fault rates of ORAM-backed paging.
}
CLOCK page replacement maintains an index of the next page to be replaced (i.e., the \emph{clock hand}).
When subpage replacement is requested, \thename\ checks the access bit of the corresponding \gls{pte} of the block pointed to by the current hand.
If the access bit is set, \thename\ clear the bit, increment the clock hand and repeat the previous step until a non-access page is encountered.
When a non-access page is reached, the page pointed to by the clock hand is chosen for eviction, and the clock hand is once again incremented.


\hdr{Handling page dependencies}
When dealing with subpages for data and page tables, \thename\ must consider the dependencies between these subpages before selecting pages to be swapped out of the scratchpad area.
To maintain the \gls{gva}-to-\gls{gpa} translation chain, \thename\ must not swap out page table pages that currently have active mappings.
Hence, before selecting a page for eviction, \thename\ scans the \glspl{pte} of the pages to locate any present \glspl{pte}.
If a present \gls{pte} is found, it follows the \gls{pfn} encoded in the first \gls{pte} to access the \gls{gpt} subpage of the lower level within the same page.
This process repeats until a data subpage is encountered (reaching page level 0) or no more present \glspl{pte} are found at the current \gls{gpt} level.
The last visited subpage is then selected for eviction, and the clock hand is incremented, following the CLOCK replacement algorithm.

% When dealing with both subpages for data and page tables, \thename must also consider the \emph{subpage dependencies}, before choosing pages to be swapped out of the scratchpad area.
% That is, \thename\ must not swap out page table pages that currently have active mapping, or the \gls{gva}-to-\gls{gpa} translation chain would be broken.
%
% Hence, before selecting a page for eviction, it scans the \glspl{pte} of the pages to locate any present \glspl{pte}. If a present \gls{pte} is found, it follows the \gls{pfn} encoded in the \emph{first} \gls{pte} to access the \gls{gpt} subpage of the lower level within the same page.
% This dependency check resembles a software page table walk, with the distinction that it visits the first present \gls{pte} found. This process repeats until a data subpage is encountered (reaching page level 0), or no more present \glspl{pte} are found at the current \gls{gpt} level.
% The last visited subpage is then selected for eviction, and the clock hand is incremented, following CLOCK replacement.

% While CLOCK replacement is sufficient when only data subpages are considered,

% \thename instead chooses to preemptively \emph{prevent} the \gls{gva}-to-\gls{gpa} translation chain from being broken by subpage replacement.
% It introduces a dependency-aware subpage replacement method built upon the CLOCK replacement algorithm.
% Instead of immediately evicting a page when the clock hand points to it, \thename checks for dependencies by inspecting its \glspl{pte}.

% Finally, \thenaem\ aim to preserve the layout of the unikernel application.
% Previous attack techniques leverage the controlled channel to infer
% First, Wilke et al. \cite{wilke2020sevurity} leverage the controlled channel and the deterministic page access pattern at boot time to identify
%
% Second,  \cite{morbitzer2021severity}
% Dependencies between subpages are demonstrated in \autoref{fig:replacement}.
% % For a \gls{gva}-to-\gls{gpa} translation to be active, all of the subpage \gls{gpt} used for its translation must be activated (fetched into the \gls{asa} by the \gls{pager}).
% When a subpage that has active mappings is evicted by the replacement algorithm, the chain of \gls{gva}-to-\gls{gpa} is broken.
% Thus, the \texttt{pt\_vaddr} for all of its dependent subpages (A and B in \autoref{fig:replacement}) becomes invalid, and these subpages become \emph{orphan} subpages.

% The problem is apparent later when the orphan page A is chosen for replacement; since \texttt{pt\_vaddr} is invalid, we cannot use it to unmap the subpage.
% A naive approach to solve this would be to evict all of the dependent subpages along the \gls{gva}-to-\gls{gpa} translation chain, whether a \gls{gpt} page is chosen for eviction.
% However, this leads to a surge in the number of subpages inside the ORAM stash after a subpage replacement.
% This is undesirable in ORAM schemes due to the problem of ORAM stash overflow that happens when stash occupancy is too high~\cite{stefanov2018path,freecursiveoram}.
% Another approach could be to trigger a software page table walk (\cref{sec:pf-handler}) with the ORAM block ID of A as the faulting virtual address to re-fetch the depended \gls{gpt} subpage into the \gls{asa}.
% This, however, would incur additional subpage replacements that lead to additional software page table walks.





% In Unikraft, the unikernel \thename is prototyped upon, and the virtual memory subsystem is implemented through in a \gls{vma} abstraction library called \texttt{ukvmem}.
% To support \thename, we introduce a non-intrusive extension to \texttt{ukvmem} to ensure \gls{pager}, implemented as a \gls{vma}'s \texttt{fault} callback, is invoked when handling \gls{gpt} pages during the software page table walk, similar to how Unikraft's paging handles last-level \glspl{pte}.
% \thename then implements its paging scheme as a new type of \gls{vma} called the \gls{ovma}.
% \key
% {
% 	The abstraction allows \thename to manage the unikernel and its hosted program memory using its oblivious paging scheme.
% 	It also allows \emph{exclusion} of certain pages from the paging scheme.
% 	For instance, the \gls{gpa} of an I/O buffer shared with the hypervisor is inevitably known to the hypervisor.
% }
%
% More importantly, it allows the selective treatment of the interrupt handlers, which must be excluded from \thename ORAM-backed paging to avoid \emph{double-faulting}.



% \key{
% 	\thename provides a deterministic security guarantee: in cases of page smashing attacks, the leaked information from \gls{npf} attacks on the \gls{apage} would only extract at most a configurable $n$ number of memory accesses.
% }
% Note that $n$ can be configured to  1,  at which point the validation scheme prevents the hypervisor from leveraging page smashing to leak any memory accesses.

% One concern with this approach is that the hypervisor can manipulate the value returned by \texttt{rdtsc} to affect $T$.


% In more detail, $T_{smash}$ includes the overheads taken to (1) invoke \texttt{PSMASH} on a page, (2) unmapping of the smashed \gls{gpa} page range (i.e., \emph{zapping} the page range in the \gls{npt}).
% Assuming that all \gls{cvm}'s \gls{gpa} have been unmapped in the \gls{npt} as a preparation, $T_{npf}$ consists of the overheads to (3) wait for the guest to access memory and trigger an \gls{npf} that would incur a \texttt{VMEXIT}, (4) remap the faulting \gls{gpa}  (5) invoking \texttt{VMRESUME} to resume \gls{cvm}'s execution.
% In \cref{sec:eval:psmash-profile}, we provide the detailed measurement of $T_{smash}$ and $T_{npf}$ on a \gls{sevsnp}-enabled machine.


\section{Adaptive Defense Layer}
The purpose of \thename's adaptive defense layer is to introduce an intolerable cost to the data extraction primitives in Phase II explained in ~\cref{sec:modeling-attacks}.

\subsection{Securing periodic security routines}
\thename's \gls{pao} scheme we discussed in \cref{sec:design:mmu} requires a security invariant that the hypervisor must not be able to smash the huge mapping during \gls{cvm} execution.
This requires a mechanism to \emph{periodically} check and verify the security invariants of the running \gls{cvm}.
Unfortunately, within a \gls{cvm}, there is a lack of reliable ways to intercept the running application and perform the security routines periodically.
While interrupt controllers like \gls{apic} and \gls{pit} could be programmed by the \gls{cvm} to trigger interrupt at a high, the actual interrupt delivery is virtualized by the hypervisor, making hardware-based interrupts untrustworthy.
% For instance, a malicious hypervisor can intentionally delay or not deliver virtualized interrupts, causing the security routines to be delayed or indefinitely postponed.
Previously, for Intel SGX enclaves, systems for detecting abnormal execution delays~\cite{dejavu} and enclave exit tracking~\cite{oleksenko2018varys} were proposed.
These works leverage a custom compiler to insert periodic security-checking after a few hundred instructions, hence having the source code requirement limitations we discussed in \cref{sec:preliminaries}.

\hdr{Time-verified interrupt}
% To achieve its goal of being compatible with unmodified application binaries and fully leveraging the underexplored opportunities in \glspl{cvm}, 
To ensure that security routines are invoked periodically and timely, \thename\ introduces a new OS primitive called the \emph{time-verified interrupt} (\veriint).
\veriint\ is essentially a wrapper for interrupt handler that measures the reliability of hardware interrupt delivery through the secure timer now available in \gls{cvm} (\othree).
The user of \veriint\ specifies a security callback $Cb_{periodic}$ that is expected to be invoked at each period $T$.
A small constant $\Delta$ is also specified to consider the variability of interrupt deliveries.
The actual interrupt period programmed into hardware is $T - \Delta$, so actual interrupt delivery falls into $T \pm \Delta$.
The user also registers $Cb_{late}$ that is to be invoked whenever the interrupt delivery is detected to be unreliable, which may perform \gls{cvm} termination or rerandomization of scratchpad (\cref{sec:design:rerandomization}).

\veriint\ guarantees the reliability of interrupt injection through two thesholds, $\delta_{late}$ and $\delta_{reset}$.
The \veriint\ interrupt handler increments a \emph{lateness} counter if an interrupt is detected late, and $Cb_{late}$ is invoked whether the counter surpasses $\delta_{late}$.
The lateness counter is reset every $\delta_{reset}$.
In other words, the rate of on-time interrupt must be at least $\delta_{late}/\delta_{reset}$.

\begin{algorithm}[h]
	\small
	% \SetKwFunction{Fn}{VerifIntHandler}
	\SetKwProg{Fn}{fn}{:}{}
	\SetKwFunction{FMain}{VeriIntHandler}
	\Fn{\FMain{$T$, $\Delta$, $\delta_{late}$, $\delta_{reset}$ ,$Cb_{periodic}$, $Cb_{late}$}}{
		$T_{now} \gets SecureRdtsc()$\;
		$T_{elapsed} \gets T_{now} - T_{prev}$\;
		\If {$T_{elapsed} > (T + \Delta)$}{
			$Cnt_{Late} \gets Cnt_{Late} + 1$\;
		}
		\If{$Cnt_{late} > \delta_{late}$}{
			$Cb_{Late}$()\;
		}
		\If{$Cnt_{Handled} > \delta_{reset}$}
		{
			$Cnt_{Late} \gets 0$\;
			$Cnt_{Handled} \gets 0$\;
		}
		$Cb_{periodic}$()\;
		$T_{prev} \gets SecureRdtsc()$\;
		$Cnt_{handled} \gets Cnt_{handled}+1$\;
	}
	\caption{Time-verified interrupt handler}
\end{algorithm}



\hdr{Detection of interrupt starving}
% In OSes that support preemptive scheduling, hardware timer interrupts are fundamental for OS scheduling, so postponing interrupt delivery leads to DoS.
% However, unikernels often leverage simpler cooperative scheduling that does not require hardware periodic interrupts~\cite{mirage,unikraft}, detecting the lack of interrupts must be in place.
We choose to detect long delays of interrupts \emph{opportunitstically} whenever the control is passed back from the \gls{cvm} application to the unikernel, which occurs on two occasions.
The first occasion is when handling any types of interrupts and traps, including the page faults by our \thename.
The second occasion is whenever a system call is invoked.
On any of these occasions, for each \veriint\ interrupt handler, we compute the time since the last interrupt and invoke $Cb_{late}$ if a long delay is detected.



\hdr{Page Smashing detection with \veriint}
\todo{Rewrite to reflect TOCTOU scenario}
To determine the optimal revalidation period $T$, we carefully analyze the overheads required to \emph{smash} a page and perform subsequence \gls{npf} attacks.
We first define the total time it takes to perform the page-smashing attack to extract $n$ memory accesses as  $T_{smash} + T_{npf} \times n$, where $T_{smash}$ is the overhead of smashing the \gls{apage} into 512 4KB mappings, and $T_{npf}$ is the time it takes to extract a single \gls{gpa} access, using the \gls{npf} control-channel attack.
Thus, by configuring the validation thread to re-validate the \gls{apage} at a period $T < T_{smash} + T_{npf} \times n$, the information leak obtained through page smashing would be bounded to $n$ memory accesses.




\subsection{Adaptive scratchpad rerandomization}
\thename\ adaptively rerandomize the physical layout of the scratchpads regions in response to high exit rates.

\hdr{VMEXIT tracking with \veriint}
Leveraging the verified interrupt.
\todo{wip}
Secret extraction through fine-grained primitives requires abnormally high exit rates that can be leveraged to indicate attacks~\cite{oleksenko2018varys,zhang2020klotski}.
We leverage \veriint\ to implement a periodic VMEXIT tracking scheme.
by polling for the exit code inside \gls{vmsa}.
$T$ is set to be the overhead of a round trip of VM enter and exit.
% A common characteristic of fine-grained attack primitives is that they require abnormally high rates of interruption~\cite{zhang2020klotski,oleksenko2018varys}, which is also the case for fine-grained extraction attacks targeting SEV \glspl{cvm}~\cite{li2021cipherleaks,sevstep}.
% Thus, previous works in Intel SGX leverage interrupt rate as a metric for attack detection and abort the execution the interrupt rate surpasses a chosen threshold~\cite{oleksenko2018varys,dejavu}.
% Unfortunately, detection-based approaches suffer from false positives when the application requires frequent interrupts (e.g., I/O intensive applications)~\cite{oleksenko2018varys}.

% To get away with false positive rates, \thename\ instead chooses to build a layer-2 defense on top of layer 1's \gls{pao} protection.
% It introduces a \emph{subpage randomization} scheme that is adaptive to the rate of \gls{cvm} interruption (VMEXIT).
% It first implements a VMEXIT detection scheme that leverages the secure timer (\textbf{O3}) to periodically pool for the hardware-maintained exit code inside the \gls{vmsa}.
% After the number of VMEXIT reaches a configured threshold, \thename\ \emph{rerandomize} the subpages within the \glspl{apage}, halting in-flight attack attempts.



\subsection{Selective ciphertext protection}
To cope with layer-2 rerandomization, attackers now must leverage the low-bandwidth fine-grained side channels to extract the memory accesses from the entire guest physical spaces used for \glspl{apage} (up to 4MB).
For cache attacks, this is rather infeasible because of their susceptibility to noise.
On the other hand, we chose to provide a defense against ciphertext side-channel due to their \emph{deterministic} nature that can be used to break layer 2's probabilistic guarantees.
We leverage MMU access (\textbf{O1}) to build a protection scheme against the \emph{colision}-based and \emph{dictionary}-based ciphertext side channels~\cite{cipherh,li2022systematic}.
% Against the dictionary-based ciphertext side channel, which aims to build plaintext-ciphertext dictionaries from memory content, we leverage page permissions to implement a paging-based trap-and-salt mechanism to add nonces, making every memory written to selected pages indistinguishable.
% Collision-based ciphertext side-channel observes whether a memory location changed or not.
% Against this side channel, we introduce a protection that opportunistically performs dummy memory writes to memory locations within \gls{apage}.
% Kernel-performed dummy memory writes are indistinguishable from application memory writes, thanks to the unikernel design.

\hdr{Paging-powered trap-and-emulate}

\hdr{Dummy accesses}

\hdr{Selective enforcement}


